% Documento de arquitectura de KARSWCOL (Release 2).
% Compilar en Overleaf (pdfLaTeX) o con: latexmk -pdf main.tex
% Los recuadros "PENDIENTE DEL EQUIPO" marcan datos que solo el equipo tiene
% (sprints, burndown, ceremonias): completarlos antes de entregar.
\documentclass[11pt,a4paper]{article}

\usepackage[utf8]{inputenc}
\usepackage[T1]{fontenc}
\usepackage[spanish,es-tabla]{babel}
\usepackage{lmodern}
\usepackage[margin=2.3cm]{geometry}
\usepackage{graphicx}
\usepackage{booktabs}
\usepackage{tabularx}
\usepackage{longtable}
\usepackage{array}
\usepackage{xcolor}
\usepackage{tcolorbox}
\usepackage{listings}
\usepackage{tikz}
\usetikzlibrary{arrows.meta,positioning,shapes.geometric,fit,backgrounds,calc,babel}
% Separador horizontal dentro de los nodos de clase/entidad (hline no funciona en TikZ).
\newcommand{\sep}{\\[-1pt]\rule{3.2cm}{0.4pt}\\}
\usepackage[hidelinks]{hyperref}

\definecolor{karsRed}{HTML}{CE1126}
\definecolor{karsYellow}{HTML}{FCD116}
\definecolor{karsBlue}{HTML}{003893}

\newtcolorbox{pendiente}{colback=yellow!12,colframe=orange!70!black,title=PENDIENTE DEL EQUIPO,fonttitle=\bfseries\small}
\newcolumntype{L}{>{\raggedright\arraybackslash}X}

\lstset{basicstyle=\ttfamily\small,breaklines=true,frame=single,columns=fullflexible}

\tikzset{
  comp/.style={draw,rounded corners,fill=karsBlue!8,minimum width=3.2cm,minimum height=0.9cm,align=center,font=\small},
  ext/.style={draw,dashed,rounded corners,fill=gray!8,minimum width=3cm,minimum height=0.9cm,align=center,font=\small},
  node0/.style={draw,thick,rounded corners,fill=karsYellow!15,inner sep=8pt},
  arr/.style={-{Stealth[length=2.2mm]},thick}
}

\title{\textbf{KARSWCOL}\\[4pt]\large Carreras multijugador en tiempo real con identidad colombiana\\[2pt]\normalsize Documento de arquitectura --- Release 2}
\author{Julián Tinjacá \and Sara Arteaga \and Johan Sebastián Beltrán\\[4pt]
\small Arquitecturas de Software (ARSW) --- Escuela Colombiana de Ingeniería Julio Garavito}
\date{Noviembre de 2026}

\begin{document}
\maketitle
\begin{center}
\small Juego: \url{https://karswcol.onrender.com} \quad
Repositorio: \url{https://github.com/Julianizz18/karswcol}
\end{center}
\tableofcontents
\newpage

% =====================================================================
\section{Introducción}
\subsection{Objetivo}
Describir la arquitectura de KARSWCOL, un juego de carreras multijugador en
tiempo real para navegador, y justificar cómo atiende los atributos de
calidad del Release 2 (rendimiento, escalabilidad, seguridad, disponibilidad,
portabilidad, mantenibilidad y observabilidad), con evidencia ejecutable.

\subsection{Alcance}
Cubre el cliente web, el servidor autoritativo (Node.js), el despliegue en
Render, la capa de observabilidad (Prometheus, Grafana, Loki), las pruebas
automáticas, las pruebas de carga con k6 y el pipeline de CI/CD. Las
decisiones que no se implementaron por la restricción de costo del equipo
(despliegue distribuido con balanceador) se documentan como arquitectura
objetivo y se distinguen explícitamente de lo implementado.

\subsection{Audiencia}
Docente y monitores de ARSW, y desarrolladores que quieran operar o
extender el sistema.

% =====================================================================
\section{Visión general}
\subsection{Descripción del sistema}
Hasta ocho jugadores (completados con bots si se desea) corren en circuitos
inspirados en trazados de Fórmula 1 y ambientados en cinco regiones de
Colombia (Caribe, Andes, Eje Cafetero, Llanos y Amazonía), con vehículos
icónicos (TransMilenio, taxi, burro, avión de Avianca, metro de Medellín),
objetos de la cocina colombiana (tinto, arepa, aguapanela, lechona, tamal y
buñuelo) y el clima real de cada región. El servidor es la única fuente de
verdad: simula la carrera a 20~Hz y difunde el estado por WebSocket; el
cliente solo envía intención y dibuja (2D o 3D con Three.js).

\subsection{Requerimientos funcionales clave}
\begin{itemize}
  \item HU01 --- Iniciar sesión con Google, elegir vehículo y entrar a la sala.
  \item HU02 --- Correr con física autoritativa (agarre, derrape, choques, vueltas por puntos de control).
  \item HU03 --- Recoger objetos compartidos con resolución de conflictos.
  \item HU04 --- Eventos dinámicos: aguacero, pista mojada, charcos y aquaplaning.
  \item HU05 --- Resultados con posición, tiempo total y mejor vuelta.
  \item HU06 --- Reconexión y continuidad de la sesión.
  \item Release 2 --- Anfitrión de sala, pausa compartida con límite, votación para volver al lobby, espera con entrada automática, calidad gráfica adaptativa.
\end{itemize}

\subsection{Requerimientos no funcionales clave}
\begin{tabularx}{\textwidth}{lL}
\toprule
\textbf{Atributo} & \textbf{Meta}\\\midrule
Rendimiento & Estado cada 50~ms; gap entre estados p95 $<$ 100~ms; HTTP p95 $<$ 300~ms; errores $<$ 1~\%\\
Seguridad & Sesión firmada y revocable; autorización por rol (anfitrión); límite de mensajes\\
Disponibilidad & Detección de caída en $\le$ 30~s; reconexión automática\\
Mantenibilidad & Cobertura $\ge$ 80~\% en backend y frontend; Quality Gate en verde\\
Portabilidad & Misma imagen Docker en nube y on-premise, sin cambios de código\\
Observabilidad & Métricas, logs, dashboards y alertas reproducibles\\
\bottomrule
\end{tabularx}

% =====================================================================
\section{Marco metodológico}
\subsection{Metodología}
Scrum, con Azure DevOps como herramienta oficial (Product Backlog, Sprint
Backlog, tablero y Burndown Chart). Release 1: sprints 1--5; Release 2:
sprints 6--10.

\subsection{Criterios DoR / DoD}
\textbf{Definition of Ready}: historia con criterios de aceptación
verificables, estimada en puntos, sin dependencias bloqueantes y entendida
por el equipo.\\
\textbf{Definition of Done}: código en \texttt{master} con su prueba
automática; pipeline de CI en verde (pruebas, cobertura, Quality Gate,
imagen Docker); desplegado en Render; criterios de aceptación verificados en
el juego; documentación actualizada.

\subsection{Planificación y estimación}
\begin{pendiente}
Completar con los datos de Azure DevOps: historias por sprint, puntos
estimados, capacidad (capacity) del equipo, velocidad por sprint y captura
del Burndown Chart de cada sprint del Release 2.
\end{pendiente}

% =====================================================================
\section{Atributos de calidad}
Cada escenario sigue la plantilla fuente / estímulo / artefacto / entorno /
respuesta / medida. El detalle completo y la lista de pruebas está en la
carpeta \texttt{docs/} del repositorio.

\subsection{Rendimiento}
\begin{tabularx}{\textwidth}{lL}
\toprule
Fuente / estímulo & 40 jugadores y espectadores conectados más 5 visitantes por segundo cargando el cliente\\
Artefacto / entorno & Gateway y game loop, operación normal\\
Respuesta & El servidor sostiene el tick de 50~ms y atiende HTTP y WebSocket sin errores\\
Medida & HTTP p95 $<$ 300~ms; gap entre estados p95 $<$ 100~ms; errores $<$ 1~\%\\
\bottomrule
\end{tabularx}
Resultados de carga, estrés y picos en la sección~\ref{sec:resultados-rendimiento}.

\subsection{Escalabilidad y disponibilidad}
\begin{tabularx}{\textwidth}{lL}
\toprule
Fuente / estímulo & La instancia se cae o se reinicia durante una carrera\\
Artefacto / entorno & Servicio \texttt{karswcol} en Render\\
Respuesta & El cliente avisa y reconecta con espera creciente; al volver recupera su kart (30~s de gracia); Prometheus alerta \texttt{KarswcolCaido}\\
Medida & Detección $\le$ 30~s; reconexión $\le$ 15~s tras volver el servidor; 0 acciones manuales\\
\bottomrule
\end{tabularx}

\paragraph{SLA y cálculo.} Render no publica SLA para planes distintos de
Enterprise; se usa un supuesto conservador de 0{,}995. Para el login de
Google se toma como referencia el SLA de Cloud Identity (0{,}999). Open-Meteo
no entra al cálculo porque tiene respaldo (clima típico por pista).

\begin{center}
\begin{tabular}{lccc}
\toprule
Despliegue & $A_{total}$ & Downtime mensual & Downtime anual\\\midrule
Actual (serie): $0{,}995 \times 0{,}999$ & 99{,}40~\% & 4~h 19~min & 52~h 31~min\\
Objetivo: LB $\times\,[1-(1-0{,}995)^2] \times 0{,}999$ & 99{,}89~\% & 48{,}6~min & 9~h 51~min\\
\bottomrule
\end{tabular}
\end{center}
El despliegue actual alcanza ``dos nueves''; con dos instancias en paralelo
el sistema se acerca a tres nueves y el factor limitante pasa a ser el login
de Google. Puntos únicos de falla actuales: la instancia de Render, el estado
en memoria y el proveedor de identidad para nuevos inicios de sesión. El
cálculo se contrasta con la disponibilidad observada en Grafana
(\texttt{avg\_over\_time(up[\$\_\_range])}) y con la tasa de errores de k6.

\subsection{Seguridad}
Tres escenarios (detalle en \texttt{docs/seguridad.md}):
\begin{description}
  \item[S1 Autenticación y sesión.] Una cookie alterada, vencida o de una sesión cerrada se rechaza (401 / WebSocket 4401); el cierre de sesión revoca la sesión en el servidor y desconecta sus pestañas. Medida: 100~\% de rechazos, revocación inmediata.
  \item[S2 Autorización.] Solo quien entró a la sala la controla, solo el anfitrión la configura y una cuenta tiene un solo kart. Medida: 0 cambios de estado provocados por actores sin permiso.
  \item[S3 Protección frente a ataques.] Un cliente que inunda el WebSocket ve descartado el exceso (30 msg/s, ráfaga 60) y es desconectado a los 300 descartes; la entrada mal formada se ignora; Prometheus alerta. Medida: el servicio sigue atendiendo a los demás.
\end{description}

\subsection{Portabilidad}
Escenario de migración Cloud $\rightarrow$ On-Premise: la misma imagen
(\texttt{node:22-alpine}, sin root, con \texttt{HEALTHCHECK}) corre en Render
y en un servidor propio con \texttt{docker compose up --build}, incluida la
observabilidad. Medida: $\le$ 30~min y 0 cambios de código.

\subsection{Mantenibilidad}
Escenario de inspección continua: en cada push, GitHub Actions (reusable
workflows) ejecuta las pruebas con cobertura, SonarCloud analiza el código y
espera el Quality Gate, y se construye y prueba la imagen Docker.

\begin{center}
\begin{tabular}{lccc}
\toprule
Capa & Líneas & Ramas & Meta\\\midrule
Backend (\texttt{server/src}) & 95{,}1~\% & 82{,}9~\% & $\ge$ 80~\%\\
Frontend (\texttt{client/js}) & 95{,}1~\% & 84{,}9~\% & $\ge$ 80~\%\\
\bottomrule
\end{tabular}
\end{center}

\subsection{Observabilidad}
El gateway expone \texttt{/metrics} (Prometheus) y \texttt{/health}; los logs
son JSON estructurados recogidos por Promtail hacia Loki. Grafana se
aprovisiona con los dashboards ``KARSWCOL --- Servidor'' y ``KARSWCOL ---
Pruebas de carga (k6)''. Condiciones de degradación con alerta: servidor
caído (\texttt{up == 0}), tick p95 $>$ 35~ms, intervalo entre ticks p95 $>$
75~ms, lag del event loop p99 $>$ 100~ms, memoria $>$ 400~MB, errores 5xx, y
alertas de seguridad (mensajes descartados, conexiones rechazadas, inicios
de sesión fallidos).

% =====================================================================
\section{Vista de arquitectura}
\subsection{Estilo arquitectónico}
Cliente--servidor con \textbf{servidor autoritativo} y comunicación
\textbf{orientada a eventos} por WebSocket. El backend se organiza en
componentes de responsabilidad única (Gateway, Matchmaking, Game Engine,
Concurrency Resolver, Event Manager, Ranking, Weather, Observabilidad)
orquestados por la sala (\texttt{RaceRoom}).

\subsection{Modelo cloud}
PaaS (Render, Web Service con Docker) para el juego; servicios SaaS
externos para identidad (Google) y clima (Open-Meteo); GitHub Actions para
CI y SonarCloud para análisis estático.

\subsection{Decisiones clave}
\begin{tabularx}{\textwidth}{lL}
\toprule
\textbf{Decisión} & \textbf{Justificación}\\\midrule
Servidor autoritativo a 20~Hz & Evita trampas y da un orden total de eventos por tick, que resuelve la concurrencia sin bloqueos\\
WebSocket con difusión por tick & Latencia baja y estado consistente para todos; el cliente interpola a 60~fps\\
Sesión firmada (HMAC) en cookie HttpOnly & El WebSocket se autentica con la misma cookie; sin estado en el servidor salvo la lista de revocación\\
Una sala por proceso & El game loop necesita memoria local y latencia mínima; se escala agregando instancias que alojan salas\\
Render gratuito & Restricción de costo del equipo; se documenta la arquitectura objetivo distribuida\\
\bottomrule
\end{tabularx}

\subsection{Balanceo, despliegue distribuido y tiempo real}
La arquitectura objetivo coloca un balanceador con afinidad de sesión
(\emph{sticky}) delante de varias instancias del servidor, de modo que el
WebSocket de un jugador permanezca en la instancia que aloja su sala. Redis
(ya orquestado) mantiene el directorio de salas, las sesiones revocadas y el
canal Pub/Sub entre instancias; PostgreSQL (ya orquestado) guarda historial
y ranking. Por la restricción de costo, el despliegue implementado usa una
sola instancia en Render; la justificación cuantitativa está en el cálculo
de disponibilidad.

% =====================================================================
\section{Diagramas UML}
\subsection{Contexto}
\begin{center}
\begin{tikzpicture}[node distance=1.2cm and 1.6cm]
  \node[comp,fill=karsYellow!25] (sys) {KARSWCOL};
  \node[ext,left=of sys] (player) {Jugador\\(navegador)};
  \node[ext,above right=0.6cm and 1.6cm of sys] (google) {Google Identity};
  \node[ext,right=of sys] (meteo) {Open-Meteo};
  \node[ext,below right=0.6cm and 1.6cm of sys] (ops) {Operador\\(Grafana)};
  \draw[arr] (player) -- node[above,font=\scriptsize]{HTTPS / WSS} (sys);
  \draw[arr] (sys) -- node[sloped,above,font=\scriptsize]{verifica ID token} (google);
  \draw[arr] (sys) -- node[above,font=\scriptsize]{clima} (meteo);
  \draw[arr] (ops) -- node[sloped,below,font=\scriptsize]{métricas y logs} (sys);
\end{tikzpicture}
\end{center}

\subsection{Componentes}
\begin{center}
\begin{tikzpicture}[node distance=0.55cm and 0.5cm]
  \node[comp] (gw) {Connection Gateway\\\scriptsize HTTP, WebSocket, /metrics};
  \node[comp,below=of gw] (room) {RaceRoom\\\scriptsize orquestador de sala};
  \node[comp,left=of room,xshift=-0.4cm] (auth) {Auth Service};
  \node[comp,below left=0.7cm and -0.6cm of room] (mm) {Matchmaking};
  \node[comp,below=0.7cm of room] (ge) {Game Engine};
  \node[comp,below right=0.7cm and -0.6cm of room] (cr) {Concurrency\\Resolver};
  \node[comp,below=of mm] (em) {Event Manager};
  \node[comp,below=of ge] (rk) {Ranking};
  \node[comp,below=of cr] (bot) {Bot Driver};
  \node[comp,right=of room,xshift=0.4cm] (wx) {Weather Service};
  \node[comp,above=of gw] (client) {Game Client\\\scriptsize app.js, race3d.js};
  \draw[arr] (client) -- node[right,font=\scriptsize]{WebSocket} (gw);
  \draw[arr] (gw) -- (room);
  \draw[arr] (gw) -| (auth);
  \foreach \t in {mm,ge,cr,em,rk,bot} \draw[arr] (room) -- (\t);
  \draw[arr] (room) -- (wx);
\end{tikzpicture}
\end{center}

\subsection{Clases (núcleo del servidor)}
\begin{center}
\begin{tikzpicture}[node distance=0.6cm and 0.8cm, every node/.style={font=\scriptsize}]
  \node[draw,align=left,fill=white] (room) {\textbf{RaceRoom}\sep state, trackId, hostId\\ playerStates, disconnected\sep join(), handleDisconnect()\\ setReady(), togglePause()\\ startRace(), tick(), finishRace()};
  \node[draw,align=left,fill=white,below left=of room] (ge) {\textbf{GameEngine}\sep createPlayerState()\\ setInput(), step()\\ evaluateCheckpoint()};
  \node[draw,align=left,fill=white,below=of room] (cr) {\textbf{ConcurrencyResolver}\sep pickups[respawnAt]\sep tryClaim(), serialize()};
  \node[draw,align=left,fill=white,below right=of room] (em) {\textbf{EventManager}\sep rainLap, rainDurationMs\sep onLapCompleted(), tick()};
  \node[draw,align=left,fill=white,left=of room] (mm) {\textbf{MatchmakingService}\sep players\sep join(), hasCapacity()\\ allReadyToStart()};
  \node[draw,align=left,fill=white,right=of room] (rk) {\textbf{RankingService}\sep buildRanking()\\ persistRace()};
  \foreach \t in {ge,cr,em,mm,rk} \draw[-{Diamond[open]},thick] (\t) -- (room);
\end{tikzpicture}
\end{center}

\subsection{Entidad--Relación (modelo objetivo en PostgreSQL)}
El estado de juego vive en memoria; el modelo persistente objetivo para el
ranking es:
\begin{center}
\begin{tikzpicture}[node distance=2.2cm, every node/.style={font=\scriptsize}]
  \node[draw,align=left] (u) {\textbf{jugador}\sep \underline{sub} (Google)\\ nombre, foto};
  \node[draw,align=left,right=of u] (r) {\textbf{resultado}\sep \underline{carrera\_id}, \underline{sub}\\ posición, tiempo\_ms, mejor\_vuelta\_ms};
  \node[draw,align=left,right=of r] (c) {\textbf{carrera}\sep \underline{id}\\ pista, dificultad, clima, fecha};
  \draw (u) -- node[above]{1 \quad N} (r);
  \draw (r) -- node[above]{N \quad 1} (c);
\end{tikzpicture}
\end{center}

\subsection{Secuencia: entrar a la sala y recoger un objeto}
\begin{center}
\begin{tikzpicture}[every node/.style={font=\scriptsize}]
  \foreach \x/\n in {0/Cliente,3.5/Gateway,7/RaceRoom,10.5/ConcurrencyResolver} {
    \node[draw,fill=karsBlue!8] at (\x,0) {\n};
    \draw[dashed] (\x,-0.3) -- (\x,-7.2);
  }
  \draw[arr] (0,-0.8) -- node[above]{WS + cookie de sesión} (3.5,-0.8);
  \draw[arr] (0,-1.5) -- node[above]{join} (3.5,-1.5);
  \draw[arr] (3.5,-2.0) -- node[above]{join(conn, cuenta)} (7,-2.0);
  \draw[arr,dashed] (7,-2.5) -- node[above]{lobby\_update (a todos)} (0,-2.5);
  \draw[arr,dashed] (3.5,-3.0) -- node[above]{joined (id)} (0,-3.0);
  \draw[arr] (0,-3.7) -- node[above]{input \{up,left\}} (3.5,-3.7);
  \draw[arr] (3.5,-4.2) -- node[above]{setInput} (7,-4.2);
  \node[draw,fill=yellow!20] at (7,-4.9) {tick (50 ms)};
  \draw[arr] (7,-5.5) -- node[above]{tryClaim(jugador)} (10.5,-5.5);
  \draw[arr,dashed] (10.5,-6.0) -- node[above]{pickup\_claimed} (7,-6.0);
  \draw[arr,dashed] (7,-6.7) -- node[above]{event + state (a todos)} (0,-6.7);
\end{tikzpicture}
\end{center}

\subsection{Despliegue}
\begin{center}
\begin{tikzpicture}[node distance=0.8cm and 1.2cm]
  \node[comp] (browser) {Navegador\\\scriptsize cliente web (Three.js)};
  \node[comp,right=2.2cm of browser] (svc) {Contenedor karswcol\\\scriptsize Node 22 (gateway + game loop + cliente estático)};
  \begin{scope}[on background layer]
    \node[node0,fit=(svc),label={[font=\scriptsize]above:{Render (PaaS, plan gratuito, HTTPS)}}] {};
  \end{scope}
  \node[ext,below=1.3cm of svc] (google) {Google Identity};
  \node[ext,right=0.6cm of google] (meteo) {Open-Meteo};
  \node[ext,below=1.3cm of browser] (gh) {GitHub Actions\\\scriptsize CI + SonarCloud};
  \draw[arr] (browser) -- node[above,font=\scriptsize]{HTTPS / WSS} (svc);
  \draw[arr] (svc) -- (google);
  \draw[arr] (svc) -- (meteo);
  \draw[arr] (gh) -- node[sloped,below,font=\scriptsize]{push $\rightarrow$ redeploy} (svc);
\end{tikzpicture}
\end{center}
Entorno de pruebas reproducible (un solo host, Docker Compose): gateway,
Redis, PostgreSQL, Prometheus, Grafana, Loki, Promtail y k6.

\subsection{Actividades: ciclo de una carrera}
\begin{center}
\begin{tikzpicture}[node distance=0.45cm, every node/.style={font=\scriptsize}]
  \node[draw,rounded corners] (a) {Lobby: entrar, elegir (anfitrión), listo};
  \node[draw,rounded corners,below=of a] (b) {Cuenta regresiva 3-2-1};
  \node[draw,rounded corners,below=of b] (c) {Carrera: tick 50 ms (input, física, objetos, lluvia, estado)};
  \node[draw,diamond,aspect=3,below=of c] (d) {¿todos llegaron o se agotó el tiempo?};
  \node[draw,rounded corners,below=of d] (e) {Resultados y votación para volver al lobby};
  \draw[arr] (a) -- (b); \draw[arr] (b) -- (c); \draw[arr] (c) -- (d);
  \draw[arr] (d) -- node[right]{sí} (e);
  \draw[arr] (d.east) -- ++(1.2,0) |- node[near start,right]{no} (c.east);
  \draw[arr] (e.west) -- ++(-1.4,0) |- (a.west);
\end{tikzpicture}
\end{center}

% =====================================================================
\section{Detalles tecnológicos}
\begin{tabularx}{\textwidth}{lL}
\toprule
Lenguajes y frameworks & JavaScript (ES2022); Node.js 22 con \texttt{ws}, \texttt{google-auth-library}, \texttt{prom-client}; cliente sin framework con Three.js para la vista 3D\\
Infraestructura & Render (Docker, HTTPS), Google Identity Services, Open-Meteo\\
DevOps y CI/CD & GitHub Actions con reusable workflows (pruebas y cobertura, SonarCloud con Quality Gate, imagen Docker); despliegue automático en Render; Azure DevOps para Scrum\\
Docker / Compose & \texttt{server/Dockerfile} (Node 22, sin root, HEALTHCHECK); \texttt{docker-compose.yml} con gateway, Redis, PostgreSQL y observabilidad\\
Observabilidad & Prometheus (métricas y alertas), Grafana (dashboards aprovisionados), Loki y Promtail (logs)\\
Pruebas & \texttt{node:test}; jsdom para el cliente; k6 para carga, estrés y picos\\
\bottomrule
\end{tabularx}

% =====================================================================
\section{Seguridad del sistema}
\subsection{Autenticación y autorización}
Inicio de sesión con Google Identity Services: el servidor verifica el ID
token (firma, audiencia y vencimiento) y emite una cookie de sesión propia
firmada con HMAC-SHA256 (\texttt{HttpOnly}, \texttt{SameSite=Lax},
\texttt{Secure} bajo HTTPS, 12~h deslizantes). Cada sesión tiene un
identificador revocable. La autorización se aplica en el servidor: conexión
del mismo origen con sesión válida, control de la sala solo tras entrar,
configuración solo para el anfitrión, un kart por cuenta.

\subsection{Cifrado}
TLS terminado por Render (HTTPS y WSS). Las cookies se firman (integridad);
no se almacenan datos personales más allá de la sesión.

\subsection{Buenas prácticas}
Sin secretos en el repositorio; contenedor sin root; validación de entrada;
límite de tamaño de cuerpo HTTP (16~KB) y de mensajes WebSocket; métricas
protegidas con token en despliegues públicos; endpoint de pruebas de carga
deshabilitado en producción.

\subsection{Evidencias de los tres escenarios}
Pruebas en \texttt{server/test/auth.test.js}, \texttt{room.test.js},
\texttt{input.test.js}, \texttt{metrics-auth.test.js} y
\texttt{client-lobby.test.js}; alertas validadas con \texttt{promtool}.
\begin{pendiente}
Adjuntar la aprobación del profesor de los tres escenarios y capturas de
las alertas de seguridad en Prometheus/Grafana.
\end{pendiente}

% =====================================================================
\section{Bitácora de ceremonias}
\begin{pendiente}
Sprint 0; Release 1 (sprints 1--5) y Release 2 (sprints 6--10): fecha,
objetivo y resultado de cada Planning, Review y Retrospectiva, con enlaces
o capturas de Azure DevOps. Lecciones aprendidas del equipo.
\end{pendiente}

% =====================================================================
\section{Conclusiones y recomendaciones}
\subsection{Logros}
Juego multijugador en tiempo real desplegado y accesible; servidor
autoritativo con resolución determinista de concurrencia probada con
clientes simultáneos; cobertura del 95~\% en ambas capas; pipeline con
análisis estático; observabilidad completa con alertas; pruebas de carga,
estrés y picos con resultados reproducibles.

\subsection{Riesgos}
Instancia única sin SLA (plan gratuito) y arranque en frío; estado en
memoria; dependencia de Google para nuevos inicios de sesión.

\subsection{Mejora continua}
Implementar la arquitectura objetivo (balanceador, varias instancias,
Redis y PostgreSQL); monitor externo de disponibilidad; persistencia del
ranking.

\subsection{Resultados de rendimiento, disponibilidad y observabilidad}
\label{sec:resultados-rendimiento}
Ver \texttt{docs/rendimiento.md} (tablas generadas con
\texttt{load-tests/resumen.mjs} a partir de los resúmenes de k6) y
\texttt{docs/disponibilidad.md}.
\begin{pendiente}
Copiar aquí las tablas finales de \texttt{docs/rendimiento.md} y las
capturas de Grafana de cada prueba (antes, durante y después).
\end{pendiente}

% =====================================================================
\section{Anexos}
\subsection*{Glosario}
\begin{description}
  \item[Tick] Ciclo fijo de simulación del servidor (50~ms).
  \item[Gap entre estados] Tiempo entre dos mensajes de estado consecutivos en el cliente.
  \item[Anfitrión] Jugador con más antigüedad en la sala; configura pista, hora, dificultad y bots.
  \item[TPS / TPM] Transacciones por segundo / por minuto.
\end{description}
\subsection*{Repositorios y evidencias}
\begin{itemize}
  \item GitHub: \url{https://github.com/Julianizz18/karswcol} (CI en \emph{Actions})
  \item Azure DevOps: \href{https://dev.azure.com/kartombia-arsw/Proyecto\%20ARSW\%202026-2}{dev.azure.com/kartombia-arsw (Proyecto ARSW 2026-2)}
  \item Scripts k6: \texttt{load-tests/}; resultados: \texttt{docs/evidencias/k6/}
  \item Dashboards: \texttt{observability/grafana/dashboards/}; alertas: \texttt{observability/prometheus/alerts.yml}
\end{itemize}
\subsection*{Referencias}
\begin{enumerate}
  \item L. Bass, P. Clements y R. Kazman, \emph{Software Architecture in Practice}, 4.\textsuperscript{a} ed. Addison-Wesley, 2021.
  \item Google Cloud, ``Identity Platform Service Level Agreement''. \url{https://cloud.google.com/identity-platform/sla}
  \item Google Cloud, ``Cloud Identity Service Level Agreement''. \url{https://cloud.google.com/terms/identity/sla}
  \item Render, ``Status''. \url{https://status.render.com}
  \item Grafana Labs, ``k6 documentation''. \url{https://grafana.com/docs/k6/latest/}
  \item Prometheus, ``Alerting rules''. \url{https://prometheus.io/docs/prometheus/latest/configuration/alerting_rules/}
\end{enumerate}

\end{document}
